\SourceLeaf{298}
From $\dfrac{6^\ast\times 8^\ast}{11^\ast}$ or $\dfrac{6^\ast\times 10^\ast}{9^\ast}$ or $\dfrac{8^\ast\times 10^\ast}{7^\ast}$ it follows that

\[
PP=qq,
\]

from which one must necessarily conclude that

\[
P=+q
\]

provided one considers only triangles whose sides and angles do not exceed $180^0$. Then one has immediately

\[
Q=+s,\qquad R=+p,\qquad S=+r,
\]

which together are the four desired propositions.

I leave it to you to decide how much use to make in your work of this or of my other customary ways of arranging calculations. \mbox{.\,.\,.}

\medskip

\Needspace{8\baselineskip}
\begin{center}
[2.]

\textsc{Gauss} to \textsc{Gerling}. Göttingen, 26 June 1816.
\end{center}

\mbox{.\,.\,.\,.\,.\,.\,.} I have been so constantly distracted by various matters that I have not yet got around to considering the most suitable order for the first propositions of spherical trigonometry in a geometrical treatment. I should think, however, that it would be simplest to take Euclid XI. 20. 21 as the basis and to derive the proposition that $A+B+C>180^0$ from the polar triangle. I scarcely think that any substantial advantage can be gained by considering the plane triangle formed by the vertices of the spherical one, since there do not seem to be any \textls{very} simple relations among its angles. The following may well be one of the most elegant. If $A'$, $B'$, $C'$ are the angles of the plane triangle and

\[
A+B+C-180^0=u,
\]

then

\[
\cos A'=\cos\tfrac{1}{2}a\cos(A-\tfrac{1}{2}u).
\]

\mbox{.\,.\,.\,.\,.\,.\,.\,.\,.}
